% Title page: a strong typographic hierarchy, the hook, and the boxed abstract.
\begin{titlepage}
\thispagestyle{empty}
\setstretch{1.0}
\setlength{\parindent}{0pt}
\vspace*{0.1in}

{\sffamily\footnotesize\color{stone}\textls[160]{\MakeUppercase{Working Paper}}
 \hfill \textls[60]{\MakeUppercase{\paperseries}}\par}
\vspace{6pt}
{\color{indigo}\rule{\linewidth}{1.2pt}}
\vspace{0.42in}

{\fontsize{34}{40}\selectfont\color{indigo}\raggedright
 Less Is a Different\\[2pt] Kind of More\par}
\vspace{14pt}
{\fontsize{17}{21}\selectfont\itshape\color{ink!85}\papersubtitle\par}
\vspace{0.3in}

{\sffamily\normalsize\color{ink}\paperauthor\par}
\vspace{2pt}
{\sffamily\small\color{stone}\paperdate\par}

\vspace{0.32in}
\begin{pullquote}[fontupper=\Large\itshape\color{indigo}]
The light bulb needed a power grid.\\
The bottle light needed a roof and the sun.
\end{pullquote}
\vspace{0.2in}

\begin{abstractbox}
\setlength{\parindent}{0pt}
When people lack the resources that an existing solution takes for granted,
they are sometimes forced to rethink the problem, and the result can be
cheaper, simpler and more widely useful than the original. This paper argues
that scarcity does not produce more innovation but, under particular
conditions, a different kind. I set out a three-part mechanism (removal of
the default, reframing and reverse flow) and test it against off-grid solar
and mobile money in sub-Saharan Africa, low-cost health care in India, research
on creativity under constraints, and historical cases of wartime and island
substitution, using public data throughout. The evidence supports a narrower
claim than the one I began with. Mobile money now reaches more Kenyan adults
than banks do, but the number of people in sub-Saharan Africa without
electricity is higher than it was a quarter-century ago, and the frugal
solutions that do work rest on capabilities built by abundance. I conclude that
abundance creates capabilities and scarcity sometimes remixes them into forms
that reach more people, when basic tools, room to experiment and a path to
scale are present. Far more often, scarcity simply means going without.
\end{abstractbox}

\vfill
{\sffamily\scriptsize\color{stone}
 Every number in this paper is generated from the data pipeline in the
 accompanying repository. Working paper: comments and corrections are welcome.\par}
\end{titlepage}
